%----------
%	ARQUITECTURA Y DISEÑO AGÉNTICO
%----------

El núcleo de esta investigación compara dos formas de poner la metodología experta a
disposición del modelo. Ambas resuelven la misma tarea, sobre el mismo corpus y
con el mismo número de llamadas al modelo. Lo único que cambia entre ellas es
cómo llega esa metodología.

\begin{itemize}
    \item \textbf{Nivel B0 (\textit{baseline}).} La metodología de cada variable
    se inyecta entera en el \textit{prompt}. El modelo emite su veredicto en una
    sola llamada, sin herramientas.
    \item \textbf{Nivel B1 (\textit{Agent Skills}).} El modelo recibe únicamente
    los metadatos de las \textit{skills} disponibles (su nombre y una
    descripción) y carga bajo demanda el cuerpo de la que necesita. Este mecanismo
    se conoce como \textit{progressive disclosure}. Además puede consultar las
    guías expertas y verificar sus evidencias mediante herramientas.
\end{itemize}

\begin{figure}[H]
\centering
\includegraphics[width=0.94\textwidth]{arquitectura_b0_b1_blanco.pdf}
\caption{Los dos niveles comparados. Ambos hacen cinco llamadas por artículo (una
por variable) y comparten formato de salida.}
\label{fig:iris-b0b1}
\end{figure}

El diseño fija todo salvo un factor: ambos niveles hacen \textbf{cinco llamadas
por artículo}, una por variable, y comparten el mismo formato de salida. Lo único
que cambia es cómo llega la metodología al modelo. De este modo, el delta de
acuerdo entre B0 y B1 es atribuible al empaquetado como \textit{Agent Skills}, y no
al número de llamadas ni al formato de la respuesta.

\subsection{Diseño agéntico: un agente por dentro}
\label{subsec:diseno_agentico}

Antes de describir cómo se organiza la arquitectura  multi-agéntica, conviene fijar de qué se
compone \emph{uno}. Para ello hay que distinguir dos conceptos que a menudo se
confunden. Una \textit{skill}\footnote{El concepto de \textit{Agent Skills} fue
formalizado por Anthropic a finales de 2025 como un estándar abierto para la
portabilidad entre plataformas: \url{https://www.anthropic.com/engineering/equipping-agents-for-the-real-world-with-agent-skills} [Último acceso: 05/08/2026]}
es conocimiento empaquetado: un fichero de metodología, pasivo, que se carga pero
no decide. Un \textbf{agente} es un ejecutor con un bucle de decisión, que elige
qué skills cargar y cuándo verificar. \textit{Una skill por variable no es un agente} pero
\textit{un agente que carga su skill bajo demanda} sí lo es, y es lo que aquí se
implementa.

Cada agente dispone de un \textbf{catálogo de skills} del que, de entrada, solo ve
el nombre y una descripción; carga el cuerpo de una skill únicamente cuando la
necesita. Este mecanismo se conoce como \textit{progressive disclosure}, un
principio del diseño de interacción persona-máquina (HCI por sus siglas en inglés), mostrar solo lo esencial y revelar el detalle
bajo demanda, que se popularizó en el ámbito de la
usabilidad~\cite{nielsen1994} y que la comunidad de agentes ha tomado
prestado para el diseño del contexto. La idea es la misma: cargar únicamente lo
relevante para la tarea en curso evita saturar la ventana de contexto con
información que no viene al caso, algo que degrada la respuesta y encarece la
inferencia. No procede de una fuente académica primaria sobre agentes, sino de esa
tradición de usabilidad; su formulación operativa para modelos es la del marco de
\textit{Agent Skills} citado más arriba. El catálogo reúne tres tipos de skill: la
de la propia variable (obligatoria), unas auxiliares compartidas y unos resúmenes
de las guías expertas, cuyo formato y generación se detallan en la
Sección~\ref{prompts}.

El agente opera mediante un \textbf{bucle estilo ReAct}~\cite{yao2023react} (de
\textit{Reasoning} y \textit{Acting}), un patrón que entrelaza razonamiento y
acción en vez de pedir la respuesta de un solo tirón. En cada turno el modelo razona sobre qué le falta,
emite \emph{una} acción, observa el resultado que le devuelve el sistema y vuelve a
razonar con esa nueva información; el ciclo se repite hasta que cierra con su
veredicto. Ese entrelazado es lo que le permite, por ejemplo, cargar su skill,
consultar una guía si duda y verificar sus evidencias antes de responder, en lugar
de contestar a ciegas.

Las acciones se expresan \textbf{en texto}, no mediante el \textit{tool-calling}
nativo de cada proveedor. La razón es práctica: el enrutado multi-proveedor se
apoya en una interfaz de texto común, y el \textit{tool-calling} nativo es frágil e
incompatible entre modelos, sobre todo en los locales pequeños. Con un protocolo de
texto, el mismo mecanismo funciona igual en los cuatro proveedores. En cada turno
el agente emite una de estas acciones y el sistema le devuelve el resultado:


\begin{itemize}
    \item \texttt{LEER\_SKILL <id>}: carga el cuerpo de una skill.
    \item \texttt{CONSULTAR\_GUIA <consulta>}: recupera pasajes literales de las
    guías expertas (Sección~\ref{rag}).
    \item \texttt{VERIFICAR ["..."]}: comprueba que las evidencias son literales
    del texto.
    \item \texttt{FINAL \{...\}}: cierra con el veredicto de la variable.
\end{itemize}

El bucle admite hasta unas ocho iteraciones. Si el agente emitiera \texttt{FINAL}
sin haber usado ninguna herramienta, su comportamiento equivaldría al del
\textit{baseline}; ese «colapso a B0» se registra como una métrica más, aunque en
la práctica apenas se observa, porque todos los agentes cargan al menos su propia
skill. Más informativa resulta la \emph{tasa de uso} de cada herramienta: con qué
frecuencia un modelo consulta las guías o verifica sus evidencias. Esa medida de
comportamiento se reporta en el Capítulo~\ref{results_and_discussion} y ayuda a
interpretar por qué la arquitectura mejora el acuerdo en unos modelos y no en
otros.

\begin{figure}[H]
\centering
\includegraphics[width=\textwidth]{agente_anatomia_blanco.pdf}
\caption{Anatomía de un agente especializado, con el ejemplo del agente de
\texttt{lenguaje\_sexista}: catálogo de skills (progressive disclosure) y bucle de
acciones que carga la skill bajo demanda, consulta guías, verifica evidencias y
cierra con el veredicto.}
\label{fig:iris-agente-anatomia}
\end{figure}

\subsection{Arquitectura multiagente}
\label{subsec:multiagente}

Sobre ese diseño se construye la arquitectura del nivel B1: no un único agente que
juzgue las cinco variables a la vez, sino \textbf{cinco agentes especializados},
uno por variable (V25, V26, V30, V33 y V35). Cada uno tiene un único objetivo,
clasificar su variable, y resuelve su tarea de forma independiente de los demás,
con su propio catálogo y su propio bucle. Cada artículo se procesa, por tanto, con
cinco agentes y cinco llamadas, una por variable.

La especialización responde a una decisión de diseño. Sobrecargar a un modelo con
mucho contexto tiende a degradar su rendimiento~\cite{liu2023lost}: cuando la
información relevante se diluye entre la que no lo es, aumentan las
alucinaciones~\cite{huang2025hallucination} y la confusión entre criterios. Un solo agente encargado de las cinco variables tendría que sostener a
la vez cinco metodologías y sus fronteras, incurriendo en ese riesgo. Acotar cada
agente a una variable reduce la carga de contexto, mantiene el foco y hace que su
traza (qué skills cargó, si verificó, si colapsó a B0) sea interpretable variable a
variable.

\begin{figure}[H]
\centering
\includegraphics[width=\textwidth]{agentes_enjambre_blanco.pdf}
\caption{Arquitectura multiagente del nivel B1: cada artículo se procesa con cinco
agentes especializados independientes, uno por variable, cada uno con su propio
catálogo y su propio bucle.}
\label{fig:iris-multiagente}
\end{figure}
